% ============================================================
% Sección 6
% ============================================================
\section{6 Análisis Cualitativo (Geométrico)}

\begin{tcolorbox}[enunciado]
\textbf{Clases:}
\begin{itemize}[leftmargin=*]
    \item Isoclinas, Análisis Cualitativo 1 MCA4 2Sep26
    \item Análisis Cualitativo asíntotas y Sol Analítica, EDO Autónoma 7 Sep 2026 MCA4
\end{itemize}

\textbf{Notas:}
\begin{itemize}[leftmargin=*]
    \item \texttt{Ecuacion\_y }$y^{\prime}=g(y)$\texttt{ .2016\_17.1 Sep2024.pdf}
\end{itemize}
\end{tcolorbox}

\vspace{0.3cm}
\textbf{Responde:}

\begin{enumerate}[label=\arabic*., leftmargin=*]
    \item ¿Qué se entiende por una EDO autónoma?
    
    \item ¿Cómo son las isoclinas de una EDO de la forma:
    \[ \frac{dy}{dx}=f(x) \]
    
    \item ¿Cómo son las isoclinas de una EDO de la forma (autónoma):
    \[ \frac{dy}{dx}=f(y) \]
\end{enumerate}

\vspace{0.3cm}
\noindent \textbf{Haz el análisis geométrico (cualitativo) básico de las siguientes EDO considerando:}
\begin{itemize}[leftmargin=*]
    \item Soluciones de equilibrio (o puntos críticos)
    \item Crecimiento o decrecimiento (criterio de la primera derivada)
    \item Isoclinas y Campo Direccional
    \item Concavidad (criterio de la segunda derivada)
    \item Forma de algunas curvas integrales
\end{itemize}

\begin{enumerate}[label=\arabic*., leftmargin=*, start=4]
    \item \textbf{Logística.} Usa el ejemplo 2.1 de las notas.
    \[ P^{\prime}(t)=P\left(1-\frac{P}{2}\right) \]
    
    \item Usa el ejemplo 2.2 de las notas.
    \[ y^{\prime}=\sqrt{|y|} \]
    
    \item De los incisos anteriores, Logística y $y^{\prime}=\sqrt{|y|}$, cuál es la diferencia radical, en lo referente a su soluciones de equilibrio.
    
    \item EDO derivada del sistema Lotka-Volterra.
    \begin{equation*}
        \frac{dy}{dx}=\frac{y(x-3)}{x(5-y)} \tag{5}
    \end{equation*}
\end{enumerate}